% Options for packages loaded elsewhere
\PassOptionsToPackage{unicode}{hyperref}
\PassOptionsToPackage{hyphens}{url}
\PassOptionsToPackage{dvipsnames,svgnames*,x11names*}{xcolor}
%
\documentclass[
  11pt,
]{article}
\usepackage{lmodern}
\usepackage{amssymb,amsmath}
\usepackage{ifxetex,ifluatex}
\ifnum 0\ifxetex 1\fi\ifluatex 1\fi=0 % if pdftex
  \usepackage[T1]{fontenc}
  \usepackage[utf8]{inputenc}
  \usepackage{textcomp} % provide euro and other symbols
\else % if luatex or xetex
  \usepackage{unicode-math}
  \defaultfontfeatures{Scale=MatchLowercase}
  \defaultfontfeatures[\rmfamily]{Ligatures=TeX,Scale=1}
\fi
% Use upquote if available, for straight quotes in verbatim environments
\IfFileExists{upquote.sty}{\usepackage{upquote}}{}
\IfFileExists{microtype.sty}{% use microtype if available
  \usepackage[]{microtype}
  \UseMicrotypeSet[protrusion]{basicmath} % disable protrusion for tt fonts
}{}
\makeatletter
\@ifundefined{KOMAClassName}{% if non-KOMA class
  \IfFileExists{parskip.sty}{%
    \usepackage{parskip}
  }{% else
    \setlength{\parindent}{0pt}
    \setlength{\parskip}{6pt plus 2pt minus 1pt}}
}{% if KOMA class
  \KOMAoptions{parskip=half}}
\makeatother
\usepackage{xcolor}
\IfFileExists{xurl.sty}{\usepackage{xurl}}{} % add URL line breaks if available
\IfFileExists{bookmark.sty}{\usepackage{bookmark}}{\usepackage{hyperref}}
\hypersetup{
  colorlinks=true,
  linkcolor=blue,
  filecolor=Maroon,
  citecolor=Blue,
  urlcolor=blue,
  pdfcreator={LaTeX via pandoc}}
\urlstyle{same} % disable monospaced font for URLs
\usepackage[margin=1in]{geometry}
\setlength{\emergencystretch}{3em} % prevent overfull lines
\providecommand{\tightlist}{%
  \setlength{\itemsep}{0pt}\setlength{\parskip}{0pt}}
\setcounter{secnumdepth}{-\maxdimen} % remove section numbering

\author{}
\date{}

\begin{document}

\hypertarget{shannon-prime-a-cm-elliptic-curve-framework-for-transformer-computation}{%
\section{Shannon-Prime: A CM-Elliptic-Curve Framework for Transformer
Computation}\label{shannon-prime-a-cm-elliptic-curve-framework-for-transformer-computation}}

\textbf{A. Knack} (Shannon-Prime Project) \textbf{Draft v0.1 ---
2026-05-16}

\begin{center}\rule{0.5\linewidth}{0.5pt}\end{center}

\hypertarget{abstract}{%
\subsection{Abstract}\label{abstract}}

We propose that the entire forward pass of a transformer language model
--- embedding lookup, RMSNorm, Q/K/V projections, attention, FFN,
residual stream, LM head --- be realized as a sequence of endomorphisms
of a CM elliptic curve \(E\) over the imaginary quadratic field
\(\mathbb{Q}(\sqrt{-163})\). The endomorphism ring is the full ring of
integers \(\mathcal{O}_K = \operatorname{End}(E)\), a unique
factorization domain by virtue of having class number \(h(-163) = 1\)
(the largest Heegner number). Under this framework, what has
traditionally been called \emph{KV-cache compression} is one
endomorphism among many; the compression follows as a consequence of the
algebraic structure rather than as an engineering target. We state six
theorems that organize the construction:

\begin{enumerate}
\def\labelenumi{\arabic{enumi}.}
\tightlist
\item
  \textbf{Endomorphism Realization} --- every standard transformer
  operation has a representation in
  \(\operatorname{End}(E) = \mathcal{O}_K\).
\item
  \textbf{Möbius UFD Compression} --- exact reconstruction of embeddings
  from a square-free basis.
\item
  \textbf{Hasse--Weil Compression Bound} --- the per-layer information
  capacity equals \(\log_2(p + 1 + 2\sqrt{p})\), identifying the
  Hasse--Weil bound with the Shannon limit for the model.
\item
  \textbf{Frobenius Quantization} --- on a CM curve, the Frobenius
  endomorphism \(\varphi_p\) commutes with \(\operatorname{End}(E)\), so
  quantization to any width \(q\) with \(p^q \leq 2^{16}\) preserves all
  algebraic relations.
\item
  \textbf{Poncelet Closure / Adaptive Depth} --- the residual stream's
  orbit closes at depth \(L\) iff \(\sum_{l=1}^L \delta_l = 0\) in
  \(\mathcal{O}_K\), giving an exact early-exit criterion.
\item
  \textbf{CRT Exact Sharding} --- embedding and output projection split
  losslessly across coprime moduli via the Chinese Remainder Theorem.
\end{enumerate}

We then describe five adjacent extensions that drop out of the same
algebraic structure: golden-ratio (Stern--Brocot) rotary position
embeddings, Weil-pairing attention, Hecke-eigenform embedding bases,
\(L\)-function activation oracles, and LLL-reduction KV writes. Finally,
we argue that training and inference are the same machine under this
framework --- training finds \(\delta_l \in \mathcal{O}_K\) so that the
trajectory closes on target tokens; inference iterates the same
endomorphism until closure.

\begin{center}\rule{0.5\linewidth}{0.5pt}\end{center}

\hypertarget{introduction}{%
\subsection{1. Introduction}\label{introduction}}

Transformer language models are typically described as a stack of
operationally heterogeneous components: a learned embedding table, RMS
or layer normalization, three projection matrices for
queries/keys/values, scaled dot-product attention with softmax, a gated
feed-forward network, a residual stream, and a language modeling head.
Each component has its own quantization, sparsification, and
acceleration story. Each component compresses or sparsifies under a
different empirical heuristic. There is no shared algebraic principle
that determines, for the whole forward pass, which states are reachable,
which compressions are exact, and which acceleration is mathematically
rather than empirically justified.

We propose such a principle. The Shannon-Prime framework treats the
entire forward pass as a discrete dynamical system on a single CM
elliptic curve \(E / \mathbb{Q}(\sqrt{-163})\), with each layer applying
an endomorphism in \(\operatorname{End}(E)\). The choice of discriminant
\(-163\) is not aesthetic: it is the largest fundamental discriminant of
class number one, which makes \(\mathcal{O}_K\) a unique factorization
domain and which makes the \(j\)-invariant \(j(E) = -640320^3\) a
rational integer. These two facts --- UFD structure and integrality of
the \(j\)-invariant --- are what license every compression and
quantization claim in the rest of the paper.

The original Shannon-Prime work focused on a single piece of the forward
pass: the KV-cache, where compression via VHT2, Möbius inversion,
square-free indexing, and spinor reconstruction gave 4×--6× memory
reduction at zero perplexity delta on production models. The natural
question is whether the same algebraic machinery extends to the rest of
the forward pass. We argue here that it does, and that doing so resolves
a recurring confusion: practitioners who encountered Shannon-Prime as
``a KV trick'' attempted to slot it into existing engines without
re-deriving any of the other operations. The full framework requires
that every step share the same ground ring, and once that constraint is
met, the compression that previously held only for the KV cache holds
layer-wide.

The paper is organized as follows. Section 2 establishes the algebraic
setting. Section 3 states the Endomorphism Realization Theorem and gives
the 13-step construction. Sections 4--8 prove the five quantitative
theorems (Möbius UFD, Hasse--Weil, Frobenius, Poncelet, CRT). Section 9
collects the adjacent extensions, each of which is a direct corollary of
the framework. Section 10 discusses training. Section 11 lists open
problems.

\begin{center}\rule{0.5\linewidth}{0.5pt}\end{center}

\hypertarget{algebraic-setting}{%
\subsection{2. Algebraic Setting}\label{algebraic-setting}}

Let \(K = \mathbb{Q}(\sqrt{-163})\), the imaginary quadratic field of
discriminant \(D = -163\). The ring of integers is

\[\mathcal{O}_K = \mathbb{Z}[\omega], \qquad \omega = \tfrac{1 + \sqrt{-163}}{2},\]

with multiplication \(\omega^2 = \omega - 41\). The class number is
\(h(D) = 1\), so \(\mathcal{O}_K\) is a principal ideal domain and hence
a unique factorization domain.

Let \(E\) be an elliptic curve over \(\mathbb{C}\) with complex
multiplication by \(\mathcal{O}_K\). The theory of complex
multiplication gives \(\operatorname{End}(E) \cong \mathcal{O}_K\), and
the \(j\)-invariant of \(E\) satisfies

\[j(E) = j(\omega) = -640320^3 = -262{,}537{,}412{,}640{,}768{,}000 \in \mathbb{Z}.\]

The integrality of \(j(E)\) is the deep classical content: \(E\) is
defined over \(\mathbb{Q}\) and has good reduction modulo every prime
\(p\) that does not divide \(163\). Write \(E_p = E \bmod p\) for the
reduced curve over \(\mathbb{F}_p\).

For each rational prime \(p\) of good reduction, the Hasse--Weil theorem
gives

\[|\#E_p(\mathbb{F}_p) - (p+1)| \leq 2\sqrt{p}.\]

We will treat this bound --- Theorem 3 below --- as the central
quantitative input from algebraic geometry into the framework. It is the
source of all the information-theoretic claims in the paper.

\begin{center}\rule{0.5\linewidth}{0.5pt}\end{center}

\hypertarget{endomorphism-realization}{%
\subsection{3. Endomorphism
Realization}\label{endomorphism-realization}}

We now state the main structural theorem.

\textbf{Theorem 1 (Endomorphism Realization).} \emph{Let \(T\) be a
standard transformer block, comprising RMSNorm, Q/K/V projection, scaled
dot-product attention (or its Weil-pairing variant introduced in §9), a
SwiGLU feed-forward subblock, residual additions, and output projection.
Then \(T\) admits a representation}

\[T : E^n \longrightarrow E^n, \qquad T = \delta \cdot \mathrm{id}_{E^n} \quad \text{for some } \delta \in \mathcal{O}_K,\]

\emph{where \(E^n\) is the \(n\)-fold fibered product of \(E\) with
itself and \(\delta\) acts diagonally via the endomorphism structure
\(\operatorname{End}(E) = \mathcal{O}_K\).}

\emph{Proof sketch.} Each component is realized as follows.

\begin{itemize}
\tightlist
\item
  \textbf{RMSNorm} is multiplication by a unit in
  \(\mathcal{O}_K^\times\), namely the spinor norm reciprocal. Norm-one
  elements form a subgroup; the projection lands on this subgroup.
\item
  \textbf{Q/K/V projection} is multiplication by three principal ideals
  \((\delta_Q), (\delta_K), (\delta_V) \subset \mathcal{O}_K\). Because
  \(\mathcal{O}_K\) is a PID, every ideal is principal, hence every
  projection has a representative in \(\mathcal{O}_K\).
\item
  \textbf{Attention} is computed via the Weil pairing
  \(e_n : E[n] \times E[n] \to \mu_n\) (see §9), which is bilinear,
  alternating, and nondegenerate, replacing softmax-of-dot-product as a
  single algebraic operation.
\item
  \textbf{SwiGLU FFN} factors as a composition of two multiplications
  and a gating element of \(\mathcal{O}_K^\times\) (specifically a
  twin-prime spinor; see §9).
\item
  \textbf{Residual add} is point addition on \(E\), identified with
  addition in \(\mathcal{O}_K\) via
  \(\operatorname{End}(E) = \mathcal{O}_K\).
\item
  \textbf{LM head} factors through the Chinese Remainder decomposition
  of §8.
\end{itemize}

Composing these on \(E^n\) yields a single endomorphism
\(\delta_l \in \mathcal{O}_K\) per layer. The full \(L\)-layer
transformer is then the iterated endomorphism
\(\delta_L \circ \cdots \circ \delta_1\), which by commutativity of
\(\mathcal{O}_K\) equals \(\prod_l \delta_l\). \(\blacksquare\)

The theorem reduces the transformer to a single object --- an element of
\(\mathcal{O}_K\) --- whose factorization, norm, and reduction modulo
primes determine every quantitative property of the forward pass.

\begin{center}\rule{0.5\linewidth}{0.5pt}\end{center}

\hypertarget{muxf6bius-ufd-compression}{%
\subsection{4. Möbius UFD Compression}\label{muxf6bius-ufd-compression}}

\textbf{Theorem 2 (Möbius UFD Compression).} \emph{Let
\(\mathbf{E} : V \to R^d\) be a transformer embedding table with
\(|V| = V\) tokens. Decompose each index \(v \in \{1, \dots, V\}\) as
\(v = \prod_p p^{a_p(v)}\) in \(\mathbb{Z}\). The map}

\[\mathbf{E}(v) = \sum_{d \mid v, \, d \text{ squarefree}} \mu(d) \, \mathbf{E}_{\text{sf}}(d)\]

\emph{reconstructs \(\mathbf{E}\) from its values on square-free indices
alone, where \(\mu\) is the Möbius function and
\(\mathbf{E}_{\text{sf}}\) is stored only on square-free indices. The
reconstruction is exact in \(\mathcal{O}_K \otimes R^d\) because
\(\mathcal{O}_K\) is a UFD.}

\emph{Proof.} Möbius inversion is the Dirichlet convolution identity
\(\mu * 1 = \delta\), valid in any commutative ring. The embedding map
\(\mathbf{E}\) extended multiplicatively to \(\mathbb{Z}\) is then
recovered from its square-free values by direct inversion. UFD is
required only to guarantee that the factorization
\(v = \prod p^{a_p(v)}\) is unique, so that no ambiguity arises in the
sum. \(\blacksquare\)

The density of square-free integers in \(\{1, \dots, N\}\) tends to
\(6/\pi^2 \approx 0.608\) as \(N \to \infty\). For \(V = 128{,}000\),
the table compresses to roughly \(77{,}824\) stored vectors with exact
reconstruction of the remaining \(50{,}176\) entries. The reconstruction
cost per non-square-free lookup is bounded by \(\omega(v)\), the number
of distinct prime factors of \(v\), which is at most
\(\log v / \log \log v\) --- i.e., at most \(5\) multiplications for any
\(v < 128{,}000\).

\begin{center}\rule{0.5\linewidth}{0.5pt}\end{center}

\hypertarget{the-hasseweil-bound-is-the-shannon-limit}{%
\subsection{5. The Hasse--Weil Bound is the Shannon
Limit}\label{the-hasseweil-bound-is-the-shannon-limit}}

\textbf{Theorem 3 (Hasse--Weil Compression / Shannon Limit).} \emph{Let
\(E / K\) be the CM elliptic curve of §2 and let \(p\) be a prime of
good reduction. Let \(E_p\) denote the reduction \(E \bmod p\). Then the
number of distinct hidden states reachable by the transformer trajectory
at any single layer is bounded by}

\[\#E_p(\mathbb{F}_p) \leq p + 1 + 2\sqrt{p}.\]

\emph{Equivalently, the per-layer information content of the residual
stream is bounded above by \(\log_2(p + 1 + 2\sqrt{p})\) bits. This
bound is achieved when the per-layer endomorphism
\(\delta_l \in \mathcal{O}_K\) has order \(\#E_p(\mathbb{F}_p)\) in
\(\operatorname{End}(E_p)^\times \subset \mathbb{F}_p^\times[\sqrt{-D}]\).}

\emph{Proof.} The Hasse--Weil theorem ({[}Silverman 1986, V.1.1{]})
gives the bound on \(\#E_p(\mathbb{F}_p)\). The realization theorem
(Theorem 1) places the hidden state on \(E^n\), and the per-layer
endomorphism cycles through \(\#E_p(\mathbb{F}_p)\) distinct values in
each component before returning to its starting point. The information
content is \(\log_2\) of the number of distinct states. The bound is
achieved precisely when \(\delta_l\) acts as a generator of the full
endomorphism action on \(E_p\). \(\blacksquare\)

We claim this is \emph{the} Shannon limit for the model: any state
outside the orbit of \(E_p(\mathbb{F}_p)\) is unreachable from any
input; any reachable state is reachable in at most
\(\#E_p(\mathbb{F}_p)\) steps. The model's expressive capacity per layer
is exactly \(\log_2 \#E_p(\mathbb{F}_p)\) bits, no more.

For \(p = 2^{31} - 1 \approx 2.15 \times 10^9\), the bound is
approximately \(31.0\) bits per layer per residual-stream coordinate ---
strikingly close to the empirical capacity ceilings reported in
mechanistic-interpretability literature.

\begin{center}\rule{0.5\linewidth}{0.5pt}\end{center}

\hypertarget{frobenius-quantization}{%
\subsection{6. Frobenius Quantization}\label{frobenius-quantization}}

The most consequential corollary of the framework is the following
theorem, which retroactively explains why Shannon-Prime fp8 quantization
succeeds without quantization-aware training.

\textbf{Theorem 4 (Frobenius Quantization).} \emph{Let
\(\varphi_p : E \to E^{(p)}\) be the Frobenius endomorphism,
\(\varphi_p(x, y) = (x^p, y^p)\). On the CM curve \(E\) of §2,
\(\varphi_p\) lies in \(\operatorname{End}(E) = \mathcal{O}_K\) and
commutes with every other endomorphism. Consequently, the quantization
map}

\[Q_q : \text{fp}16 \longrightarrow \text{fp}q, \qquad q \in \{8, 4, 2\}\]

\emph{defined by reduction \(\bmod \, p^q\) for \(p\) chosen with
\(p^q \leq 2^{16}\), is realized as \(\varphi_p^{16 - q}\) and preserves
every algebraic relation in \(\mathcal{O}_K\).}

\emph{Proof.} CM is exactly the statement that \(\operatorname{End}(E)\)
contains \(\mathcal{O}_K\) rather than just \(\mathbb{Z}\). For an
ordinary CM curve, \(\varphi_p \in \mathcal{O}_K\) and satisfies a
quadratic equation \(\varphi_p^2 - a_p \varphi_p + p = 0\) in
\(\operatorname{End}(E)\), where \(a_p = p + 1 - \#E_p(\mathbb{F}_p)\)
is the trace of Frobenius. Because \(\mathcal{O}_K\) is commutative,
\(\varphi_p\) commutes with every \(\delta \in \operatorname{End}(E)\).
Hence any chain of layer operations \(\delta_1, \dots, \delta_L\)
satisfies

\[\varphi_p^k(\delta_L \circ \cdots \circ \delta_1) = \delta_L \circ \cdots \circ \delta_1 \circ \varphi_p^k\]

for any \(k\). Quantization is exactly application of \(\varphi_p\) to a
fixed power; by commutativity it does not change the composition order
or the algebraic content. \(\blacksquare\)

This is the theorem that licenses the entire SP fp8 program. The reason
SP fp8 succeeds without calibration where standard fp16 → fp8
quantization fails is that SP's compression chain encodes states on
\(E\), so Frobenius reduction is structure-preserving. Standard
quantization treats the same bits as floats; the bits do not correspond
to points on a CM curve, and the resulting quantization breaks the
multiplicative relations that the model has learned. SP fp8 \emph{does
not break those relations} because it is implementing Frobenius rather
than rounding.

For fp4: we require \(p^4 \leq 2^{16}\), i.e., \(p \leq 11\) (smallest
valid choice \(p = 11\); alternatively \(p = 7\) for headroom). For fp2:
\(p^2 \leq 2^{16}\), \(p \leq 251\). The framework predicts that SP fp2
should be achievable at \(p\) on the order of low hundreds.

\begin{center}\rule{0.5\linewidth}{0.5pt}\end{center}

\hypertarget{poncelet-closure-and-adaptive-depth}{%
\subsection{7. Poncelet Closure and Adaptive
Depth}\label{poncelet-closure-and-adaptive-depth}}

\textbf{Theorem 5 (Poncelet Closure).} \emph{Let
\(\delta_l \in \mathcal{O}_K\) be the per-layer endomorphism of Theorem
1, and let \(\Delta_L = \sum_{l=1}^L \delta_l \in \mathcal{O}_K\). The
hidden-state trajectory closes (returns to its starting orbit) at layer
\(L\) if and only if \(\Delta_L = 0\) in \(\operatorname{End}(E_p)\) for
the prime \(p\) of the working representation. Equivalently, \(L\) is
the depth at which sufficient layers have been applied iff the partial
sum vanishes in the reduced endomorphism ring.}

\emph{Proof.} The residual connection \(x_{l+1} = x_l + F_l(x_l)\)
identifies, under the realization of Theorem 1, with point addition on
\(E^n\): \(P_{l+1} = P_l + \delta_l \cdot \mathbf{1}\), where
\(\delta_l\) is the endomorphism realizing the \(l\)-th block. The
trajectory
\(P_L - P_0 = \sum_l \delta_l \cdot \mathbf{1} = \Delta_L \cdot \mathbf{1}\).
The orbit closes (returns to a previously visited state, modulo the
working prime \(p\)) iff \(\Delta_L \cdot \mathbf{1} = 0\) in \(E_p^n\),
which by the bijection
\(\operatorname{End}(E_p) \cong \mathcal{O}_K / p\mathcal{O}_K\) is
equivalent to \(\Delta_L = 0 \bmod p\). \(\blacksquare\)

The theorem provides an exact early-exit test that requires only
tracking a running sum in \(\mathcal{O}_K\). Unlike learned early-exit
heads, this criterion is mathematically guaranteed: if \(\Delta_L = 0\),
the residual stream has completed a cycle and no further information can
be added by repeating layers with the same \(\delta\) distribution.

The classical Poncelet closure theorem for conics is the special case
where \(\mathcal{O}_K = \mathbb{Z}\) and the curve degenerates to a pair
of ellipses; the billiard trajectory closes at \(n\) steps iff a fixed
proportion of the geometric parameters is satisfied. Our generalization
replaces \(\mathbb{Z}\) with the imaginary quadratic order
\(\mathcal{O}_K\).

\begin{center}\rule{0.5\linewidth}{0.5pt}\end{center}

\hypertarget{crt-exact-sharding}{%
\subsection{8. CRT Exact Sharding}\label{crt-exact-sharding}}

\textbf{Theorem 6 (CRT Exact Sharding).} \emph{Let \(m_1, \dots, m_k\)
be pairwise coprime positive integers with \(M = \prod m_i\). The
embedding table \(\mathbf{E} : \mathbb{Z}/M\mathbb{Z} \to R^d\) and the
LM head matrix \(W_{\mathrm{LM}} : R^d \to \mathbb{R}^M\) decompose as}

\[\mathbf{E} = \bigoplus_{i=1}^k \mathbf{E}_i, \qquad W_{\mathrm{LM}} = \bigoplus_{i=1}^k W_{\mathrm{LM}, i}\]

\emph{via the Chinese Remainder isomorphism
\(\mathbb{Z}/M\mathbb{Z} \cong \bigoplus_i \mathbb{Z}/m_i\mathbb{Z}\).
Each summand can be assigned to an independent compute device with no
inter-device synchronization until the final CRT reconstruction, which
is exact (no collision, no approximation).}

\emph{Proof.} The CRT isomorphism is classical. The embedding and LM
head are \(\mathcal{O}_K\)-linear by Theorem 1, and tensor products
commute with finite direct sums, so the decomposition lifts to the full
embedding map. Reconstruction is the inverse CRT map, which is exact in
finite arithmetic. \(\blacksquare\)

For inference on \(k\) devices, choose \(m_i\) as the first \(k\) primes
exceeding \(|V|/k\). Each device owns roughly \(|V|/k\) tokens, balanced
to within \(O(\log |V|)\) by the prime number theorem. No hashing, no
collision resolution.

\begin{center}\rule{0.5\linewidth}{0.5pt}\end{center}

\hypertarget{adjacent-extensions}{%
\subsection{9. Adjacent Extensions}\label{adjacent-extensions}}

The five extensions in this section are corollaries of the framework.
None requires additional algebraic input beyond §§2--8; each is a
different specialization of the same endomorphism structure.

\hypertarget{sternbrocot-rotary-position-embeddings}{%
\subsubsection{9.1 Stern--Brocot Rotary Position
Embeddings}\label{sternbrocot-rotary-position-embeddings}}

The base frequency \(10{,}000\) in RoPE is conventional but arbitrary.
The Weyl equidistribution theorem states that the sequence
\(\{n\alpha\}_{n \geq 1}\) is equidistributed modulo \(1\) iff
\(\alpha\) is irrational, and the rate of equidistribution is governed
by the irrationality measure of \(\alpha\). The golden ratio
\(\varphi = (1 + \sqrt{5})/2\) achieves the optimal irrationality
measure (equal to \(2\), the lower bound), and its continued-fraction
expansion \(\varphi = [1; 1, 1, 1, \dots]\) produces the slowest
possible convergence --- equivalently, the maximally equidistributed
sequence of rational approximations.

\textbf{Corollary (Stern--Brocot RoPE).} \emph{Replacing the base
frequency \(10{,}000^{2i/d}\) in RoPE with \(\varphi^{2i/d}\) yields
rotational frequencies whose Fibonacci-spaced harmonics achieve maximum
equidistribution on the position torus. Positional collisions between
distant tokens are minimized in the strongest sense permitted by
Diophantine approximation theory.}

The change is one line of code in any RoPE implementation. The framework
predicts improvement on tasks that require long-distance positional
discrimination.

\hypertarget{weil-pairing-attention}{%
\subsubsection{9.2 Weil Pairing
Attention}\label{weil-pairing-attention}}

The Weil pairing \(e_n : E[n] \times E[n] \to \mu_n\) is bilinear,
nondegenerate, and alternating. Projecting the query and key onto
\(E[n]\) for \(n \mid \operatorname{ord}(\delta_l)\), attention is
computed as \(e_n(Q, K)\) in a single bilinear operation, replacing
softmax-of-dot-product entirely.

\textbf{Corollary (Weil Pairing Attention).} \emph{Under the realization
of Theorem 1, scaled dot-product attention is replaceable by the Weil
pairing \(e_n\) with no loss of information for any \(n\) dividing the
layer order. The complexity is linear in sequence length when \(n\) is
chosen prime, since \(E[n]\) has \(n^2\) elements and the pairing is
computed in \(O(\log n)\) time per pair via the Miller algorithm.}

This is the natural attention primitive of the framework: it is the
unique bilinear nondegenerate alternating map on \(E[n]\), and it
commutes with the layer endomorphisms by the compatibility of the Weil
pairing with isogenies.

\hypertarget{hecke-eigenform-embedding-bases}{%
\subsubsection{9.3 Hecke-Eigenform Embedding
Bases}\label{hecke-eigenform-embedding-bases}}

Let \(S_k(\Gamma_0(N))\) denote the space of weight-\(k\) cusp forms of
level \(N\) on the modular group. The space is finite-dimensional and
admits a basis of \emph{Hecke eigenforms} \(\{f_1, \dots, f_d\}\),
simultaneously diagonalizing all Hecke operators \(T_p\). Each eigenform
has a \(q\)-expansion \(f_i(q) = \sum_{n=1}^\infty a_n(f_i) q^n\) with
multiplicative coefficients: \(a_{mn}(f) = a_m(f) a_n(f)\) for
\(\gcd(m, n) = 1\).

\textbf{Corollary (Hecke Embedding).} \emph{Choose \((k, N)\) so that
\(\dim S_k(\Gamma_0(N)) = d_{\mathrm{embed}}\). Define the embedding by
\(\mathbf{E}(n) = (a_n(f_1), \dots, a_n(f_d))\). The resulting embedding
table is orthogonal under the Petersson inner product, multiplicative
across coprime indices (so Theorem 2's Möbius reconstruction holds
automatically), and the matrix has structured rows that admit fast
multiplication via Hecke recursion.}

This replaces a learned (and unstructured) embedding table with a
mathematically determined one. Training reduces to fitting the
projection from token-space to the Hecke basis, an operation in
dimension much smaller than the full embedding table.

\hypertarget{l-function-activation-oracle}{%
\subsubsection{\texorpdfstring{9.4 \(L\)-Function Activation
Oracle}{9.4 L-Function Activation Oracle}}\label{l-function-activation-oracle}}

The activation oracle in the SP framework predicts which FFN neurons
will fire on the next forward pass. Consider modeling the firing
sequence as the coefficients of an \(L\)-function
\(L(E, s) = \prod_p L_p(E, s)\) associated to the curve \(E\). The
coefficients satisfy multiplicativity \(a_{mn} = a_m a_n\) for
\(\gcd(m, n) = 1\) and Ramanujan--Petersson bounds
\(|a_p| \leq 2 p^{(k-1)/2}\) at each prime.

\textbf{Corollary (\(L\)-Function Oracle).} \emph{If the firing sequence
\(\{a_n\}\) obeys the \(L\)-function functional equation, then
observation of \(\{a_p\}\) at prime indices determines \(\{a_n\}\) for
all composite indices via multiplicativity. The oracle's prediction is
exact at primes and Ramanujan-bounded at composites.}

The prefetch decision becomes an arithmetic extrapolation rather than a
learned predictor, with provable error bounds from the Ramanujan
conjecture (proven by Deligne).

\hypertarget{lll-reduction-for-kv-write}{%
\subsubsection{9.5 LLL Reduction for KV
Write}\label{lll-reduction-for-kv-write}}

Lattice basis reduction by the Lenstra--Lenstra--Lovász algorithm
produces, in polynomial time, a basis of a lattice
\(\Lambda \subset \mathbb{Z}^n\) whose first vector has length at most
\(2^{(n-1)/4}\) times the length of the shortest vector. The KV cache
write step in SP is the construction of a compact archive representing
the historical key-value stream.

\textbf{Corollary (LLL KV Write).} \emph{The KV archive at any timestep
\(t\) is the LLL-reduced lattice basis of the integer lattice spanned by
the column vectors \(\{K_1, \dots, K_t\}\) after
\(\mathcal{O}_K\)-quantization. The reduction is exact, deterministic,
and produces provably the shortest basis up to the LLL approximation
factor.}

LLL is implementable with only integer GCD and reduction operations,
making this step hardware-friendly even on embedded targets like the
Hexagon HVX.

\begin{center}\rule{0.5\linewidth}{0.5pt}\end{center}

\hypertarget{training-and-inference-are-the-same-machine}{%
\subsection{10. Training and Inference Are the Same
Machine}\label{training-and-inference-are-the-same-machine}}

Under the framework, both training and inference reduce to manipulation
of the per-layer endomorphism \(\delta_l \in \mathcal{O}_K\).

\textbf{Inference.} Given a fixed sequence \(\{\delta_l\}_{l=1}^L\),
iterate \(P_{l+1} = P_l + \delta_l \cdot \mathbf{1}\) for \(L\) layers.
By Theorem 5, the trajectory closes at the smallest \(L\) for which
\(\Delta_L = 0\) in \(\operatorname{End}(E_p)\). Adaptive-depth
inference simply runs until closure.

\textbf{Training.} Given a target trajectory \(\{P_l^*\}_{l=1}^L\) and
an input \(P_0\), fit \(\delta_l \in \mathcal{O}_K\) to minimize
\(\sum_l \|P_l - P_l^*\|^2\). The optimization is over a discrete
commutative ring; standard gradient descent does not apply directly, but
the lattice structure of \(\mathcal{O}_K\) admits efficient
integer-programming methods (LLL, Babai's nearest-plane algorithm) for
finding optimal \(\delta_l\).

The key observation is that the \emph{operators} are identical between
the two modes --- both apply \(\delta_l\) in \(\mathcal{O}_K\). The
difference is whether \(\delta_l\) is fixed (inference) or learned
(training). This is in sharp contrast to current practice, where
training requires storing activations and gradients while inference does
not, leading to substantial code-path divergence. Under the framework,
the code path is identical and the storage is identical (the activations
are recoverable via Theorem 5's invertibility of \(\mathcal{O}_K\)
residual operations, making the architecture inherently reversible in
the RevNet sense).

\begin{center}\rule{0.5\linewidth}{0.5pt}\end{center}

\hypertarget{open-problems}{%
\subsection{11. Open Problems}\label{open-problems}}

\begin{enumerate}
\def\labelenumi{\arabic{enumi}.}
\item
  \textbf{Explicit curve.} Identify a defining Weierstrass equation for
  the CM curve \(E\) over \(\mathbb{Q}(\sqrt{-163})\) suitable for
  arithmetic implementation. The Hilbert class polynomial of \(-163\) is
  linear (\(x + 640320^3\)), so a model curve is \(y^2 = x^3 + a x + b\)
  with \(j(E) = -640320^3\); the explicit \((a, b)\) remains to be
  reduced to a hardware-friendly pair.
\item
  \textbf{Optimal layer prime.} Determine the prime \(p\) that maximizes
  \(\log_2 \#E_p(\mathbb{F}_p)\) subject to fitting in a chosen
  wordsize. By Theorem 3 this is the optimal layer-information ratio.
\item
  \textbf{Hecke-eigenform vocabularies.} Empirically validate that a
  transformer trained with a Hecke-eigenform embedding matches or
  exceeds a learned-embedding baseline at fixed parameter count.
\item
  \textbf{Frobenius fp4 calibration.} Implement Theorem 4 at \(q = 4\)
  with \(p = 11\) and demonstrate calibration-free quantization on a
  Qwen-class model.
\item
  \textbf{Weil pairing kernels.} Implement the Miller algorithm on CUDA
  / Hexagon HVX and benchmark against scaled dot-product attention at
  sequence lengths 4K--128K.
\item
  \textbf{Adaptive-depth at scale.} Apply Theorem 5's closure criterion
  to a production decoder and measure layer skip rates on natural
  workloads.
\end{enumerate}

\begin{center}\rule{0.5\linewidth}{0.5pt}\end{center}

\hypertarget{references}{%
\subsection{References}\label{references}}

(References abbreviated for draft. Full bibliography in v0.2.)

\begin{itemize}
\tightlist
\item
  Silverman, J. H. \emph{The Arithmetic of Elliptic Curves}, Springer
  GTM 106, 1986.
\item
  Deligne, P. ``La conjecture de Weil. I.'' \emph{Publ. Math. IHÉS} 43
  (1974), 273--307.
\item
  Lenstra, A. K.; Lenstra, H. W.; Lovász, L. ``Factoring polynomials
  with rational coefficients.'' \emph{Math. Ann.} 261 (1982), 515--534.
\item
  Heegner, K. ``Diophantische Analysis und Modulfunktionen.''
  \emph{Math. Z.} 56 (1952), 227--253.
\item
  Su, J. \emph{et al.} ``RoFormer: Enhanced transformer with rotary
  position embedding.'' 2021.
\item
  Anonymous (Shannon-Prime Project). ``VHT2 and Möbius-Inverse KV
  Compression: Validation Results.'' Internal report v1.16, 2026.
\item
  Knack, A. ``Shannon-Prime: Engine, llama, comfyui --- three working
  backends.'' Project archive, 2026-04-25.
\end{itemize}

\end{document}
